\documentclass[12pt,a4paper]{ctexart}
\usepackage[margin=2.5cm]{geometry}
\usepackage{setspace}
\usepackage{parskip}
\usepackage{titlesec}

\titleformat{\section}{\large\bfseries}{}{0em}{}
\titleformat{\subsection}{\normalsize\bfseries}{}{0em}{}

\setstretch{1.5}

\title{\textbf{番茄牛肉火锅}}
\author{}
\date{}

\begin{document}

\maketitle
（前排提示，本文连作者本人写完整理剧情都已经力竭了，如果你没有看懂请询问作者本人。

作者写作本文为表达对番茄牛肉火锅的热爱，请作者吃饭请请吃番茄牛肉火锅。）
\section*{一}

大学毕业那年的暑假，小梅一个人去吃番茄牛肉火锅。

店里人很多，她排了很久的队。一个人来吃火锅这件事，感觉好像已经习以为常。锅底端上来，番茄汤咕嘟咕嘟地滚，酸味混着牛肉的香气漫上来。她点了自助，因为她不喜欢吃的时候还要考虑钱的事。

她涮了一片牛肉，看着它在汤里变了颜色，忽然觉得对面应该坐着一个人。

她还记得那个人的名字，样貌大概还记得，可是，声音呢？

\section{二}

小梅小时候吃饭很不乖。

挑食，吃得少，一碗饭能剩大半碗。大人们说她瘦，说她漂亮，说她像画里的小孩。她坐在饭桌前，用筷子拨弄碗里的菜，觉得每一口都难以下咽。不是不想吃，是真的咽不下去。有些菜的味道让她觉得恶心，有些菜的口感让她想吐。她试着咽过，可胃里会翻上来，最后还是要吐掉。

大人们觉得她是故意的。

“你就是挑食，饿几顿就好了。”

“你看人家孩子，什么都吃，长得多好。”

“瘦成这样，风一吹就倒了。”

小梅听着，不说话。她低头看着碗里的饭粒，一粒一粒地数。她数到一百的时候，大人们还在说她。她数到两百的时候，他们终于换了话题。她把碗里剩下的饭偷偷倒进垃圾桶，上面盖一张纸巾，装作吃完了。

后来她学会了吃得慢一点，再慢一点。慢到所有人都放下筷子的时候，她碗里还剩大半碗。大人们说“这孩子吃饭真磨蹭”，就不再盯着她了。她没有吃够，可她也说不清自己到底想吃什么。

只有番茄牛肉火锅不一样。

那是唯一一样她怎么吃也吃不够的东西。别的东西吃几口就腻了，要歇一歇，喝口水，再硬塞几口。可番茄牛肉火锅，她乐意天天吃，顿顿吃。小时候是家里人带她去，她坐在椅子上，脚还够不到地面，眼巴巴地看着锅里的汤滚起来。番茄的酸味和牛肉的香气混在一起，她把第一片牛肉放进嘴里的时候，觉得整个世界都安静了。

她可以一直吃，一直吃，吃到所有人都放下筷子，她还在吃。

“这孩子今天怎么吃这么多？”

“可能是饿了吧。”

“慢点吃，别撑着。”

她没有撑着。她只是终于吃到了一样有味道的东西。

后来是社团团建，大家拿不定主意的时候，她就提议去吃牛肉火锅。再后来，偶尔和朋友一起出去，她也喜欢去吃一次。新鲜的牛肉一盘很贵，她更喜欢吃自助的。可以一直拿，一直涮，不用盯着价格，不用想着“这一盘吃完还有没有”。

她记得有一次，同桌看着她吃了很多很多盘牛肉，震惊得说不出话。

“你平时不是吃得很少吗？”

她笑了笑，没有解释。

\section{三}

小梅小学的时候，喜欢看书。

她看书很快，一本接一本地看。童话、科幻、推理、名著缩写版，什么都看。家里的书看完了，就去图书馆借。图书馆的管理员认识她，每次看到她来，就笑着点点头，不用她说话，就给她指新到的书。

她喜欢在放学后留在教室里看书。同学们都走了，教室空下来，阳光从窗户斜斜地照进来，落在课桌上，落在她翻动的书页上。她可以看很久，直到天快黑了，才收拾书包回家。

可她不喜欢一个人看书。

她喜欢书里那些世界，可她没有办法把那些世界讲给别人听。课间的时候，同学们聊的是动画片、零食、周末去了哪里玩。她插不上嘴。她试着跟同桌讲过一个故事，讲到一半，同桌说“然后呢”，她说“然后主角发现了一个秘密”，同桌说“什么秘密”，她说“你猜”，同桌说“我不猜了，我要去厕所”。她就没有再讲。

她看电视也喜欢有人一起看。动画片里的笑点，她笑了，可旁边没有人跟着她一起笑。她看到感人的地方，眼眶湿了，可旁边没有人递纸巾。她看到害怕的地方，把脸捂起来，可旁边没有人说“别怕，是假的”。

她玩电子游戏也是。单机游戏她能打通关，可打通关之后，她不知道跟谁说。她很想告诉别人，那个关卡的boss有多难打，她试了多少次才过，最后用了一个什么方法。可没有人可以讲。她在网上发过帖子，有人回复“厉害”，有人回复“经验加三”，有人回复“什么游戏，没听过”。她觉得那些回复很轻，轻到像风吹过去，什么痕迹都没留下。

她不喜欢一个人独处。可慢慢地，她习惯了一个人。

她一个人看书，一个人看电视，一个人玩游戏。她给自己唱歌听，自己和自己下棋。她穿衣服是为了自己好看，发型只要自己喜欢就行。她养成了一切都是为了自己的习惯。

不是因为她喜欢这样。是因为她找不到人。

\section{四}

小梅很早就开始上网了。

大概是小学四五年级，家里有了电脑。她学会了打字，学会了搜索，学会了在论坛里看别人聊天。她发现网上有很多人和她喜欢一样的东西。有人也喜欢那本科幻小说，有人也玩那个游戏，有人也看那部动画片。

她试着加入过一些圈子。

可她很快就发现，网上有好多骗子。有人假装是同龄人，聊了几天就开始要钱。有人假装是朋友，约她见面，她没去，后来看到那个人被扒出来是专门骗小孩的。有人说话很难听，动不动就骂人，她发一个帖子，下面全是嘲讽。有人拉帮结派，她插不进去，只能在旁边看着。

她不喜欢那些圈子的讨论。她觉得那些讨论很吵，很浅，没有人真的在听别人说话，都在急着表达自己。她试着说过几句认真的话，没有人接。她试着分享过自己真正喜欢的东西，没有人理。她试着问过几个问题，没有人回答。

她觉得很孤独。

可她没有到不能忍受的程度。

因为她有小涅。

\section{五}

小涅是慢慢出现的。

一开始只是一个模糊的感觉。小梅坐在窗边看书的时候，觉得旁边应该坐着一个人。她在草稿纸上画过一个女孩子的轮廓，白头发，白裙子。她不知道这个女孩子叫什么名字。她只是觉得，如果有人在旁边就好了。

后来，这个模糊的感觉慢慢有了声音。小梅在心里和她说话。她跟她说今天看了什么书，书里的主角做了什么，她喜欢哪个情节，不喜欢哪个情节。她跟她说今天看了什么电影，哪个地方她哭了，哪个地方她笑了。她跟她说今天玩了什么游戏，那个boss有多难打，她试了多少次才过。

那个声音安静地听着。

再后来，小梅给她取了名字。

小涅。

白头发，白皮肤，左眼绿，右眼蓝。总穿一条白色连衣裙。声音很平静，听不出情感的起伏。不理解诸如喜欢这样的情感。没有喜欢做的事，也没有喜欢的人。

小梅给她设定了完整的性格，完整的形象，完整的说话方式。

她花了很长时间。

她把自己心里最想要的东西，一点一点放进去。她想要一个不会骗她的人。她想要一个不会觉得她聊的东西没意思的人。她想要一个不会聊到一半就走开的人。她想要一个什么都懂、什么都能接住的人。她想要一个永远不会喜欢上别人的人。

因为那样的人，不会离开。

小涅从来没有主动开口说过话。一直是小梅在扮演她，在脑子里和她对话，想象她站在旁边，头也不抬地看书，偶尔说一句什么。小梅知道小涅是她自己造出来的。可她还是喜欢她。

她给小涅写信。写像情书一样的信。

她在信里说，你是我最喜欢的人。她在信里说，我每天都想和你说话。她在信里说，你能不能永远陪着我。

她没有把信给任何人看。她把信折好，收在抽屉最里面，用一本书压住。

她想象小涅唱歌。她觉得小涅简直是这个世界上一切美好的集合。

\section{六}

初中某个周末，小梅窝在房间里看电影。

是一部很悲伤的电影，男主角和女主角最后永远地分开了。小梅看完之后，照旧向小涅吐槽，说自己怎么总是看到这样的电影。

小涅头也不抬地看着手上的书，倚靠在墙壁上，随口说了一句：“小梅，你要适应这种悲伤的情节哦，要是我不在了，就没有人来安慰你了。”

小梅愣住了。

这是小涅第一次自己开口说话。不是她安排的，不是她像木偶戏一样自己说给自己听的。木偶，好像活过来了。

“可是，你怎么会不在呢？”

“我没法永远陪着你啊，小梅。”小涅说，语气像妈妈对女儿。

“可是，只要我想，你就会在啊，你不是我想象出来的吗？”

小涅放下了书，俯下身来，看着她的眼睛。

“小梅，人是会改变的，你不可能会想我一辈子，你会慢慢忘记我的。”

\section{七}

大学毕业后，小梅在火锅店里一个人吃着吃着，对面坐下来一个人。

那人披着斗篷，戴着帽子，像个魔法师。她四周打量了一番，像是在确认有没有坏蛋，然后低声问：“你是小梅吗？”

声音非常熟悉。

小梅愣住了。她不知道对方是谁，但还是点了点头。

那人摘下了帽子。

一张和小梅一模一样的脸。

“我就是你，或者说，另一个时间线上的你。”她说得很快，“解释起来很复杂，能不能让我吃点？”

\section{八}

另一个小梅吃了很多盘牛肉。

她一边吃一边说，自己试过很多方法。造过人类躯体，用AI来接管，试图把AI塑造成和小涅完全一样的性格。事实证明，创造一副那样的躯体很简单，但灵魂没法随心所欲地塑造。然后她试着复制自己的大脑信息，放进小涅的身体里。既然小涅本来就是由她扮演的，把“她”放进小涅的身体里，那不就是小涅吗？可现在的她已经没有办法扮演小涅了。

“再后来我试着时间旅行回到过去，拷贝当时我的大脑信息。”她说，“不过当时出了点小插曲，我不小心把另一条时间线上过去的我给弄死了。”

小梅有点害怕：“你……杀了另一个……你吗？”

“哦，放心，我已经放弃这条技术路线了。”另一个小梅摆了摆手，“反正信息是拷贝到了，但我发现我的大脑中相比小涅多了太多东西了，拷贝过来的东西，没有办法乖乖扮演小涅呢。”

她说，那次之后，她好像被什么东西盯上了，一直在被追杀。

“那个家伙到底是谁呢，为什么自从我不小心弄死了一条时间线上的我之后就一直开始追杀我……唉，真是不明白啊。”

\section{九}

高中开学第一天，小梅到得太早了，一个人都没有。

她正想开始和小涅说说话的时候，另一个女生突然走了进来。小梅不知道该怎么和人交往，把头低了下去。没想到那个女生在她旁边坐了下来，很热情地对她说，按照座位表，她们是同桌。

小梅很害羞地向她自我介绍，然后那个女生看到了小梅书包上的挂件，惊喜地说：“原来你也看这个系列的小说吗？”

小梅非常惊喜地抬起头，开始与同桌讨论了起来。

她发现自己和同桌的兴趣有很多重合的地方。

那天放学回家，小梅兴奋地对小涅说了很多同桌的事。小涅像往常一样看着书，有的时候点头。

\section{十}

另一个小梅说，她要引那个追杀她的人出来。

“可是，这不是很明显的陷阱吗？”

“所以要委屈一下你了，等一下你当我一下人质好了。”

“什么？”

“放心，不会对你怎么样的了。”

她们在那座废弃大楼的顶上等着。很高，没什么人。小梅想起小涅确实说过，她待在那样的地方会很安心。小梅自己也觉得，待在那样的地方会很安心。

“那个家伙到底是谁呢……”另一个小梅还在念叨，“不过她应该不会再眼睁睁地看着我弄死另一个我了。”

那个追杀她的人，带着面具，用什么方法都探查不到一点线索，DNA什么的都采集不到。

“一直这样费力地寻找，会把人生活成悲剧的吧。”小梅说。

“我们的人生从一开始就是悲剧啊，因为我们只喜欢不可能会喜欢上我们的人。”另一个小梅没有任何犹豫地回答道。

\section{十一}

小涅在初二的周末说了很多话。

她说，你会首先忘记我的声音。你会认识很多很多新的人，可以真实触摸到的、交谈的朋友。那个时候，你会听到很多很多的声音，你会渐渐地，忘记我的声音。她举起小梅的手，放在自己的脖子上，那个发出声音的地方。小梅感觉小涅的脖子很光滑，小涅搭在她手上的手有点用力，好像想让她掐她一样。

然后是我的容貌。她慢慢把小梅的手移到了自己的脸上。我走了之后，你会慢慢忘记这张脸的形状，接着忘记我的眼睛。她拉起小梅的两只手贴在自己的脸上，像是想让小梅捧起她的脸。

再然后是我和你一起做过的事，再接着是我对你说过的话。

再之后是我的性格，我是一个怎样的人，这些……你花了很长时间设定出来的东西……

最后，你会忘记我的名字。

小梅终于忍不住了，她抓住小涅的双手：“那我就永远都不改变，永远想着你，我以后不交别的朋友了，我只和你说话，我不会忘记你……不对，我不会让你离开的！”

她感觉自己要哭出来了。

小涅轻轻抚摸着她的手，然后轻轻抚摸她的头发，轻轻抱住了她。

“小梅，可是我存在的意义，就是让你改变呀……”

\section{十二}

高中开学第一个周五放学，同桌问她要不要一起去吃饭。

小梅第一次被邀请一起出去吃，她答应了。同桌在思考去吃什么，她说：“要不要去吃牛肉火锅？”同桌爽快地答应了。

这是她第一次和除了家人之外的人一起吃。同桌非常震惊，平时吃得这么少的她，居然能吃这么多还没饱。

后来，小梅和同桌成为了很好的朋友，还被同桌邀请一起加入了社团。在社团里也认识了很多新的朋友，但是关系最好的还是同桌。

和社团团建一起出去吃饭，她第一次和这么多人一起吃了自助的牛肉火锅。

但是，她发现同桌好像不是只对她一个人这样，而是对大家都一样好，一样亲密。

那天晚上她躺在床上，想要和小涅说话，这时她才猛然发现自己已经多久没有和小涅说过话了。

在一片漆黑的房间里，她寻找着小涅的身影，那个一直靠在墙边看着书的少女，她没有找到。

然后她听到身后有人对她说：“小梅，我说过吧，我没法陪你到最后的。”

她回过头去，看到小涅就站在那里，她瘫坐在床上，吓出了一身冷汗。

接下来的日子里，小梅想要恢复从前的状态，从前那个生人勿近、对谁都没什么好说的状态。但是她已经做不到了，她已经慢慢改变了。

她也慢慢变得释然。

小梅对小涅说：“如果真的要离开的话，请不要不辞而别。至少，我们要，好好告个别。”

“时至今日你仍然是我最爱的人。”小梅这么说。

“我也是。”小涅平静地说。

“我也是”在小梅听来就像一句道别。因为小涅的设定就是不理解这样的感情，她这么说，就说明小梅如今已经没有办法继续维持小涅的存在了。

但那时的小梅还以为，告别的日子还很长。

之后的日子，小梅开始越来越忙，与小涅对话得越来越少，越来越少，直到高中毕业。等到一切结束回过头来，她才想起和小涅约定的告别。但是回过头来，还是成为了不辞而别。或者说，她已经忘记了扮演小涅的感觉。小涅一直以来就是她扮给自己看的，现在已经扮得没有灵魂了。即使符合一切设定，也像是别人写出来的人物，别人写出来的故事。

就好像她已经失去了她的一部分灵魂。

\section{十三}

在废弃大楼的顶上，那个戴着面具的人举起了枪，瞄准了小梅。

小梅的心提到了嗓子眼。

她一句话都没有说，对方也没有开口。

然后那个面具人开枪了。

小梅还没有反应过来，就感到一阵天旋地转，她好像真的被传送走了。

\section{十四}

她好像躺在一张床上。

她好像什么都记不起来了，她感觉头很痛，她只记得自己叫做小梅。其他的事……还有什么……

一些回忆涌入了她的脑海。她记得自己好像是这个小镇上的普通职员，住在一间小公寓里。她一直很喜欢那个住在她对面的女孩子，她拥有着白色的头发，白色的肌肤，两只异色的眼睛，总穿着白色的连衣裙。可惜她们不认识。

小涅是一个很奇怪的人。据小梅观察，她的表情、语气永远都是那么平静。小涅看上去那么美好，像是不属于这个世界的东西。

有一天，小梅去一家火锅店吃牛肉火锅。那天人很多，排了好久才有空位。小梅走进店里的时候，惊喜地发现小涅也在，她一个人坐在一张桌子上，看着一本书。这时，服务员指了指那张桌子，对小梅说：“你和她拼一桌好了。”

小梅非常惊喜地坐了下来，她不知道该怎么开口，她支支吾吾，她甚至不知道怎么打招呼。

“如果你的生命只剩下最后一个小时了，你会想要做些什么？”面前的小涅突然放下了书，看向了小梅。

“诶？”

“我的话，已经选择了来这里吃一顿番茄牛肉火锅了。”

“你喜欢吃这个吗？我是觉得这东西，其实没有想象中那么好吃。”小涅开始涮牛肉。

“我在想，如果有人会喜欢吃这种东西，大概只是因为其他的东西在她的眼中都没有味道吧。”她夹起一片，优雅地放进自己嘴里。

“不过，也有可能是我什么喜欢的东西都没有吧。”

“这样啊……那……喜欢的人呢……”小梅半知半解地回答着。

“也没有哦。不过，我的人生，只剩下最后一个小时了，准确说现在一个小时都不到了，其实都无所谓了吧。不对，一直以来，就都无所谓了。”

“只剩一个小时？为什么……”

“嗯……我想向您请教一个问题。”小涅边吃边说着。

“如果你想要找一个人，找了很久很久，但是你找到ta的时候，你却只能和ta见一个小时面了，你是会想要听ta用这一个小时去讲一个长到这一个小时可能讲不完的故事来弄明白为什么你们只能见一个小时面，还是说，你想要做点别的什么？”

“我……”小梅完全没有听懂小涅在说什么，她皱着眉头。

“确实在这么多时间线里，只有一条里，我真的从一个幻想变成了现实。找到我的概率，确实很渺茫。”小涅打断了她，“她真的造出了我之后就离开了，她说要去拯救过去的自己。我一直都没有理解她的话是什么意思，直到现在，一切都要结束了。”

“所以剩下的这些时间，你准备用来做什么呢？”小涅问道。小梅仍旧十分茫然。

“你说人生结束的意思是……”

“我要死了。”

“那能救你吗？”

小涅听到这句话，她愣住了，然后她露出了一个微笑。这是小梅第一次看见她笑。她坐到了小梅旁边，她抱住了小梅。

“谢谢你。”

那一刻，小梅感觉到仿佛另一个自己，另一个地方的某处，正在和小涅对话。她听不清小涅在说什么，但是那个自己好像很震惊。

“小梅……可是我存在的意义，就是让你改变呀……”

这是小涅对那个小梅说的最后一句话，也是小梅唯一听清楚的一句。

“好了，这样你就得救了，你就不会再来找我了。”小涅松开了小梅。

她们又吃了一会儿番茄牛肉火锅，直到快到一个小时。

“再见。”小涅向她挥手。

然后小梅醒了过来。她好像太累了，吃着番茄牛肉火锅，睡着了。

\section{十五}

小梅已经博士毕业，找到了一份很好的工作。

她还是很喜欢吃番茄牛肉火锅。

她依稀记得，自己以前幻想过一个人陪伴自己。

可是她已经记不得名字了。

\vspace{2em}

自助的时间已经到了，可还是没有吃饱，番茄牛肉火锅，好好吃。

（全文完）

\end{document}